% Diagram lima langkah konstruksi S-box AES.
% Dirujuk dari section-04.tex dengan % Diagram lima langkah konstruksi S-box AES.
% Dirujuk dari section-04.tex dengan % Diagram lima langkah konstruksi S-box AES.
% Dirujuk dari section-04.tex dengan % Diagram lima langkah konstruksi S-box AES.
% Dirujuk dari section-04.tex dengan \input{figures/konstruksi-sbox}.
\begin{figure}[htbp]
    \centering
    \begin{tikzpicture}[
        kotak/.style={draw, rounded corners=1pt, minimum height=1.5cm,
                      text width=2.35cm, align=center, font=\small, inner sep=2pt},
        panah/.style={-{Latex[length=2mm]}, thick},
    ]
        \node[kotak] (a)  at (0,0)    {byte masukan\\$a = xy$};
        \node[kotak] (inv) at (2.9,0)  {invers di\\$\text{GF}(2^8)$\\$a^{-1}$};
        \node[kotak] (bit) at (5.8,0)  {vektor bit\\$(b_0,\dots,b_7)^\top$};
        \node[kotak] (aff) at (8.7,0)  {matriks affine\\$8\times8$\\$\oplus\ \code{63}$};
        \node[kotak] (out) at (11.6,0) {byte keluaran\\$\text{S-box}(x,y)$};

        \draw[panah] (a) -- (inv);
        \draw[panah] (inv) -- (bit);
        \draw[panah] (bit) -- (aff);
        \draw[panah] (aff) -- (out);

        % Jejak contoh untuk a = 53.
        \node[font=\footnotesize, anchor=north] at (0,-1.05)  {$53$};
        \node[font=\footnotesize, anchor=north] at (2.9,-1.05) {$\code{CA}$};
        \node[font=\footnotesize, anchor=north] at (5.8,-1.05) {$\code{11001010}$};
        \node[font=\footnotesize, anchor=north] at (8.7,-1.05) {$\code{10110111}$};
        \node[font=\footnotesize, anchor=north] at (11.6,-1.05) {$\code{ED}$};

        \draw[dashed] (-1.35,-0.8) -- (12.95,-0.8);
        \node[font=\footnotesize, anchor=west] at (-1.35,-1.42) {contoh untuk $a = \code{53}$};
    \end{tikzpicture}
    \caption{Lima langkah konstruksi S-box AES. Byte masukan dibalik di $\text{GF}(2^8)$, hasilnya dinyatakan sebagai vektor bit, dikenai transformasi affine, lalu dibaca kembali sebagai byte}
    \label{fig:konstruksi-sbox}
    \sumber{Munir, \textit{IF4020 Kriptografi}, ITB; NIST, FIPS PUB 197.}
\end{figure}
.
\begin{figure}[htbp]
    \centering
    \begin{tikzpicture}[
        kotak/.style={draw, rounded corners=1pt, minimum height=1.5cm,
                      text width=2.35cm, align=center, font=\small, inner sep=2pt},
        panah/.style={-{Latex[length=2mm]}, thick},
    ]
        \node[kotak] (a)  at (0,0)    {byte masukan\\$a = xy$};
        \node[kotak] (inv) at (2.9,0)  {invers di\\$\text{GF}(2^8)$\\$a^{-1}$};
        \node[kotak] (bit) at (5.8,0)  {vektor bit\\$(b_0,\dots,b_7)^\top$};
        \node[kotak] (aff) at (8.7,0)  {matriks affine\\$8\times8$\\$\oplus\ \code{63}$};
        \node[kotak] (out) at (11.6,0) {byte keluaran\\$\text{S-box}(x,y)$};

        \draw[panah] (a) -- (inv);
        \draw[panah] (inv) -- (bit);
        \draw[panah] (bit) -- (aff);
        \draw[panah] (aff) -- (out);

        % Jejak contoh untuk a = 53.
        \node[font=\footnotesize, anchor=north] at (0,-1.05)  {$53$};
        \node[font=\footnotesize, anchor=north] at (2.9,-1.05) {$\code{CA}$};
        \node[font=\footnotesize, anchor=north] at (5.8,-1.05) {$\code{11001010}$};
        \node[font=\footnotesize, anchor=north] at (8.7,-1.05) {$\code{10110111}$};
        \node[font=\footnotesize, anchor=north] at (11.6,-1.05) {$\code{ED}$};

        \draw[dashed] (-1.35,-0.8) -- (12.95,-0.8);
        \node[font=\footnotesize, anchor=west] at (-1.35,-1.42) {contoh untuk $a = \code{53}$};
    \end{tikzpicture}
    \caption{Lima langkah konstruksi S-box AES. Byte masukan dibalik di $\text{GF}(2^8)$, hasilnya dinyatakan sebagai vektor bit, dikenai transformasi affine, lalu dibaca kembali sebagai byte}
    \label{fig:konstruksi-sbox}
    \sumber{Munir, \textit{IF4020 Kriptografi}, ITB; NIST, FIPS PUB 197.}
\end{figure}
.
\begin{figure}[htbp]
    \centering
    \begin{tikzpicture}[
        kotak/.style={draw, rounded corners=1pt, minimum height=1.5cm,
                      text width=2.35cm, align=center, font=\small, inner sep=2pt},
        panah/.style={-{Latex[length=2mm]}, thick},
    ]
        \node[kotak] (a)  at (0,0)    {byte masukan\\$a = xy$};
        \node[kotak] (inv) at (2.9,0)  {invers di\\$\text{GF}(2^8)$\\$a^{-1}$};
        \node[kotak] (bit) at (5.8,0)  {vektor bit\\$(b_0,\dots,b_7)^\top$};
        \node[kotak] (aff) at (8.7,0)  {matriks affine\\$8\times8$\\$\oplus\ \code{63}$};
        \node[kotak] (out) at (11.6,0) {byte keluaran\\$\text{S-box}(x,y)$};

        \draw[panah] (a) -- (inv);
        \draw[panah] (inv) -- (bit);
        \draw[panah] (bit) -- (aff);
        \draw[panah] (aff) -- (out);

        % Jejak contoh untuk a = 53.
        \node[font=\footnotesize, anchor=north] at (0,-1.05)  {$53$};
        \node[font=\footnotesize, anchor=north] at (2.9,-1.05) {$\code{CA}$};
        \node[font=\footnotesize, anchor=north] at (5.8,-1.05) {$\code{11001010}$};
        \node[font=\footnotesize, anchor=north] at (8.7,-1.05) {$\code{10110111}$};
        \node[font=\footnotesize, anchor=north] at (11.6,-1.05) {$\code{ED}$};

        \draw[dashed] (-1.35,-0.8) -- (12.95,-0.8);
        \node[font=\footnotesize, anchor=west] at (-1.35,-1.42) {contoh untuk $a = \code{53}$};
    \end{tikzpicture}
    \caption{Lima langkah konstruksi S-box AES. Byte masukan dibalik di $\text{GF}(2^8)$, hasilnya dinyatakan sebagai vektor bit, dikenai transformasi affine, lalu dibaca kembali sebagai byte}
    \label{fig:konstruksi-sbox}
    \sumber{Munir, \textit{IF4020 Kriptografi}, ITB; NIST, FIPS PUB 197.}
\end{figure}
.
\begin{figure}[htbp]
    \centering
    \begin{tikzpicture}[
        kotak/.style={draw, rounded corners=1pt, minimum height=1.5cm,
                      text width=2.35cm, align=center, font=\small, inner sep=2pt},
        panah/.style={-{Latex[length=2mm]}, thick},
    ]
        \node[kotak] (a)  at (0,0)    {byte masukan\\$a = xy$};
        \node[kotak] (inv) at (2.9,0)  {invers di\\$\text{GF}(2^8)$\\$a^{-1}$};
        \node[kotak] (bit) at (5.8,0)  {vektor bit\\$(b_0,\dots,b_7)^\top$};
        \node[kotak] (aff) at (8.7,0)  {matriks affine\\$8\times8$\\$\oplus\ \code{63}$};
        \node[kotak] (out) at (11.6,0) {byte keluaran\\$\text{S-box}(x,y)$};

        \draw[panah] (a) -- (inv);
        \draw[panah] (inv) -- (bit);
        \draw[panah] (bit) -- (aff);
        \draw[panah] (aff) -- (out);

        % Jejak contoh untuk a = 53.
        \node[font=\footnotesize, anchor=north] at (0,-1.05)  {$53$};
        \node[font=\footnotesize, anchor=north] at (2.9,-1.05) {$\code{CA}$};
        \node[font=\footnotesize, anchor=north] at (5.8,-1.05) {$\code{11001010}$};
        \node[font=\footnotesize, anchor=north] at (8.7,-1.05) {$\code{10110111}$};
        \node[font=\footnotesize, anchor=north] at (11.6,-1.05) {$\code{ED}$};

        \draw[dashed] (-1.35,-0.8) -- (12.95,-0.8);
        \node[font=\footnotesize, anchor=west] at (-1.35,-1.42) {contoh untuk $a = \code{53}$};
    \end{tikzpicture}
    \caption{Lima langkah konstruksi S-box AES. Byte masukan dibalik di $\text{GF}(2^8)$, hasilnya dinyatakan sebagai vektor bit, dikenai transformasi affine, lalu dibaca kembali sebagai byte}
    \label{fig:konstruksi-sbox}
    \sumber{Munir, \textit{IF4020 Kriptografi}, ITB; NIST, FIPS PUB 197.}
\end{figure}
